\documentclass[10pt,a4paper]{article}
\usepackage[margin=22mm]{geometry}
\usepackage{fontspec,unicode-math,amsmath,booktabs,longtable,array,fvextra,hyperref}
\setmainfont{STIXGeneral}\setmonofont{Menlo}\setmathfont{STIX Two Math}
\hypersetup{hidelinks}
\title{Complete required programs for the corrected quotient census}
\author{}\date{}\begin{document}\maketitle
The three GAP files reconstruct the marked group, evaluate all thirty literal source words, and exhaust every subgroup conjugacy class. The Python file reconciles the complete historical containment tables with the fresh census. Each file is included in full. The separate raw data file \texttt{source\_tables.json} contains all 145 printed rows; it is required data, not omitted code. The portable \texttt{check.py} orchestrates these programs, checks bindings, replays eight witnesses and tests corruptions. Its complete file is separately exported. No target list of semisparse groups is an enumeration input.
\section*{programs/ambient\_marking.g}
\VerbatimInput[fontsize=\scriptsize,breaklines=true,breakanywhere=true,breaksymbolleft={},breaksymbolright={},breakanywheresymbolpre={}]{programs/ambient_marking.g}
\section*{programs/source\_maximals.g}
\VerbatimInput[fontsize=\scriptsize,breaklines=true,breakanywhere=true,breaksymbolleft={},breaksymbolright={},breakanywheresymbolpre={}]{programs/source_maximals.g}
\section*{programs/quotient\_census.g}
\VerbatimInput[fontsize=\scriptsize,breaklines=true,breakanywhere=true,breaksymbolleft={},breaksymbolright={},breakanywheresymbolpre={}]{programs/quotient_census.g}
\section*{accounting.py}
\VerbatimInput[fontsize=\scriptsize,breaklines=true,breakanywhere=true,breaksymbolleft={},breaksymbolright={},breakanywheresymbolpre={}]{accounting.py}
\end{document}
